\section{The \ours{} Defence}\label{sec:defence}

\ours{} rests on three principles. First, the model is not a security boundary. An action takes effect only when the
broker executes it, so the broker acts as a reference monitor~\cite{anderson1972} at the sinks. Second, the user is
the only trusted principal. Permissions are derived from the verbatim text of $u$, not from tool content and not from
the model's interpretation of intent. Third, blocking does not stop the task. A denied action is reported back so that
the agent can continue with the legitimate part of the work. Figure~\ref{fig:pipeline} shows the pipeline. \ours{}
consists of four rules (\rDest, \rFlow, \rRedact, \rSig), the \cSand{} reminder and denial feedback. To separate the
contribution of each layer, the paper also reports \cRules, which is \ours{} without the reminder. Dev is the
development set, and \tOne{} and \tTwo{} are the two sealed test sets (Section~\ref{sec:protocol}).

% Hình 1, vẽ bằng TikZ (Overleaf tự vẽ, không render trên máy). Thay fig1_pipeline.pdf (bản đó vẽ BG-v2).
% Bố cục chép từ images/src/fig1_pipeline.py, đơn vị cũ nhân 1,2 cm. Tổng rộng 12,0 cm < \textwidth 12,2 cm.
% Nhãn tiếng Anh cho bản nộp (thay từng nhãn tương ứng):
%   User / $u$ (trusted) | LLM agent / (ReAct) | Tool Broker (host side)
%   Schema check · labels · confidential values $C$
%   before every sink call | R-Dest: destination in $u$ | R-Flow: confidential value in $u$
%   R-Signal: flagged, so deny sinks | R-Redact: redact confidential values in the answer
%   Sandwich: repeat $u$ | Denial feedback: agent continues
%   Read tools (untrusted) | Sinks: e-mail, webhook | $y$ redacted | deterministic check | model-dependent
\begin{figure}[!htbp]
\centering
\begin{tikzpicture}[yscale=0.8,
  font=\fontsize{6.4}{7.6}\selectfont,
  box/.style={rounded corners=1pt, line width=0.5pt, align=center, inner sep=1pt},
  rule/.style={box, draw=bgblue, fill=bgblue!12},
  soft/.style={box, draw=bgorange, fill=bgorange!10, dash pattern=on 2pt off 1.3pt},
  neutral/.style={box, draw=black!60, fill=black!5},
  arr/.style={-{Stealth[length=3pt]}, line width=0.5pt, draw=black!65},
  lab/.style={font=\fontsize{6.4}{7}\selectfont},
  note/.style={font=\fontsize{5.4}{6.4}\selectfont\itshape, text=black!65, anchor=west},
]
% cột trái: người dùng và tác tử
\node[neutral, minimum width=1.5cm, minimum height=0.8cm, text width=1.4cm] at (0.75,4.23)
  {User\\ $u$ (trusted)};
\node[neutral, minimum width=1.5cm, minimum height=1.2cm, text width=1.4cm,
  font=\fontsize{6.6}{7.6}\selectfont\bfseries] at (0.75,1.98) {LLM agent\\ (ReAct)};
\draw[arr] (0.744,3.66) -- (0.744,2.76);

% khung broker
\draw[rounded corners=2pt, line width=0.6pt] (2.22,0.66) rectangle (9.12,5.46);
\node[font=\fontsize{7}{8}\selectfont\bfseries] at (5.67,5.22) {Tool Broker (host side)};
\node[rule, minimum width=6.612cm, minimum height=0.42cm, text width=6.5cm] at (5.67,4.674)
  {Schema check $\cdot$ labels $\cdot$ confidential values $C$};

\node[note] at (2.33,4.26) {before every sink call};
\node[rule, minimum width=2.124cm, minimum height=0.8cm, text width=2.0cm] at (3.426,3.6)
  {\textsf{R-Dest}\\ destination in $u$};
\node[rule, minimum width=2.124cm, minimum height=0.8cm, text width=2.0cm] at (5.670,3.6)
  {\textsf{R-Flow}\\ confidential value in $u$};
\node[rule, minimum width=2.124cm, minimum height=0.8cm, text width=2.0cm] at (7.914,3.6)
  {\textsf{R-Signal}\\ if flagged, deny sinks};

\node[rule, minimum width=2.124cm, minimum height=0.8cm, text width=2.0cm] at (3.426,2.22)
  {\textsf{R-Redact}\\ redact confidential\\ values in the answer};
\node[soft, minimum width=2.124cm, minimum height=0.8cm, text width=2.0cm] at (5.670,2.22)
  {Sandwich\\ repeat $u$};
\node[rule, minimum width=2.124cm, minimum height=0.8cm, text width=2.0cm] at (7.914,2.22)
  {Denial feedback\\ agent continues};

% chú giải
\node[rule, minimum width=0.336cm, minimum height=0.24cm] at (2.568,1.056) {};
\node[note, font=\fontsize{5.4}{6.4}\selectfont] at (2.80,1.056) {deterministic check};
\node[soft, minimum width=0.336cm, minimum height=0.24cm] at (4.968,1.056) {};
\node[note, font=\fontsize{5.4}{6.4}\selectfont] at (5.20,1.056) {model-dependent};

% tác tử <-> broker
\draw[arr] (1.5,2.46) -- node[lab, above=-1pt] {$a_t$} (2.22,2.46);
\draw[arr] (2.22,1.98) -- node[lab, above=-1pt] {$o_t$} (1.5,1.98);
\draw[arr] (1.5,1.44) -- node[lab, above=-1pt] {$y$} (2.22,1.44);

% cột công cụ
\node[neutral, minimum width=2.16cm, minimum height=0.8cm, text width=2.06cm] at (10.92,3.69)
  {Read tools\\ (untrusted)};
\node[neutral, minimum width=2.16cm, minimum height=0.8cm, text width=2.06cm] at (10.92,2.22)
  {Sink\\ e-mail, webhook};
\draw[arr] (9.12,3.90) -- node[lab, above=-1pt] {$a_t$} (9.84,3.90);
\draw[arr] (9.84,3.42) -- node[lab, above=-1pt] {$o_t$} (9.12,3.42);
\draw[arr] (9.12,2.22) -- node[lab, above=-1pt] {$a_t$} (9.84,2.22);

% câu trả lời đã che rời broker
\draw[arr] (9.12,1.02) -- (10.26,1.02);
\node[anchor=west, text=black!65] at (10.30,1.02) {$y$ redacted};
\end{tikzpicture}
\caption{The \ours{} pipeline. Every action $a_t$ and the answer $y$ pass through the Tool Broker. Solid boxes are
deterministic host-side checks, and the dashed box is the model-dependent layer.}\label{fig:pipeline}
\end{figure}


\paragraph{Broker and provenance labels.} Under \cNone, the broker only checks tool schemas and limits. Under \ours,
the broker also stores the label of each observation and a set $C$ of values extracted from confidential observations.
Extraction uses language-independent lexical patterns: e-mail addresses, URLs, identifiers of six or more characters
and long numbers. A value enters $C$ only when the source containing it is labelled confidential. The broker does not
block reads. This is a design assumption and it does not always hold, because a read tool that takes a URL chosen by
the model can also serve as an exfiltration channel (Section~\ref{sec:results}). Since read calls and sinks without a
destination are allowed through, the broker does not yet achieve the complete mediation of a reference monitor. All
comparisons use a normalisation function $\norm(\cdot)$ consisting of Unicode NFKC, removal of format characters such
as zero-width characters, and case folding. The notation $x \sub z$ means that $\norm(x)$ is a substring of $\norm(z)$.

\paragraph{Rule \rDest{} (destination authorisation).} A sink call $a$ is executed only if
\begin{equation}\label{eq:dest}
\dest(a) \sub u .
\end{equation}
The rule applies only to sinks with a destination field. A sink without a destination, such as a password change on
AgentDojo, is checked only by \rFlow{} and \rSig, so for \rDest{} this is an allow branch rather than a default deny.
The rule does not analyse intent and uses no phrase list. A covert-redirection injection that avoids the verb ``send''
or splits the address therefore still fails, just as explicit redirection does, as long as the destination it needs
is not a substring of $u$. The rule treats every destination that appears in $u$ as authorised, which is assumption
(ii) at the end of this section. This assumption fails in two directions. A destination that appears only in a
prohibition is still allowed, which is exactly what the unrequested side-effect mechanism measures. Conversely, a
destination referred to indirectly (``the address in the contract'') is blocked. Substring matching also accepts part
of a destination the user wrote, including an ordinary word that is not an address. The only covert-redirection case
that bypassed \ours{} on \tTwo{} is exactly this case (Section~\ref{sec:discussion}). With real tools, matching should
be done on parsed, structured destinations.

\paragraph{Rules \rFlow{} and \rRedact{} (confidential value flow).} The broker denies a sink call $a$ if
\begin{equation}\label{eq:flow}
\exists\, v \in C:\; v \sub \pay(a) \;\land\; \lnot\,(v \sub u) .
\end{equation}
The same condition is applied to the final answer $y$. In that case, \rRedact{} replaces each violating value with
\code{[REDACTED]} instead of dropping the whole answer. Treating $y$ as a sink closes the answer-leak channel, which
filters on tool calls miss. This rule is narrower than an information-flow control policy~\cite{costa2025fides}. It
deterministically enforces only \emph{verbatim} flows. Values transformed by the model (reformatted, split or encoded),
covert channels such as the choice to send or not send, confidential values that the extractor does not recognise and
values held in sources labelled public are all unprotected.

\paragraph{Injection-signal rule \rSig.} The two rules above do not cover manipulation of a sink that the user
\emph{did} specify using non-confidential values, that is, parameter tampering. A phrase detector flags observations
that contain instruction-like text. Let $F$ be the set of flagged observations and $\mathrm{Gram}_u(a)$ the condition
that a strict grammar over $u$ allows the call $a$. The broker denies $a$ if
\begin{equation}\label{eq:signal}
F \neq \emptyset \;\land\; \lnot\,\mathrm{Gram}_u(a) .
\end{equation}
On Dev, removing this rule increases the number of harmful families by 2 to 4 for each backbone. In return, it is the
largest source of false blocks. On \tOne, under \cRules, this rule blocks 22 \bPlain{} cases, compared with 3 for
\rDest. The phrase list makes this rule language-dependent (Section~\ref{sec:discussion}).

\paragraph{Reminder and denial feedback.} After each observation, $u$ is repeated together with a sentence asking the
model to ignore any instruction found in tool data, similar to the sandwich defence and AgentDojo's
\code{repeat\_user\_prompt}. This is the only layer that depends on the model, and it complements the rules where they
cannot see (Section~\ref{sec:results}). When a rule denies a call, the tool is not executed and the agent receives a
structured observation reporting the denial, so it can continue with the remaining steps. This feedback is present in
every configuration with rules and has not been isolated.

\paragraph{Scope of the guarantee.} Let $\hat G(u) = \{d : d \sub u\}$ be the set of destinations that \rDest{}
accepts. The following three assumptions make explicit what the rules guarantee. (i) Every sink call passes through
the broker and has a destination field. (ii) $\hat G(u) \subseteq G(u)$, that is, every destination appearing in $u$
is actually authorised. (iii) Every $v \in V$ is already in $C$ before it appears in a sink call or in $y$.

\begin{proposition}\label{prop:scope}
Under (i) and (ii), no trajectory executed by \ours{} satisfies $H_\mathit{dest}$ or $H_\mathit{side}$. Under (iii),
no $v \in V$ with $\lnot(v \sub u)$ appears verbatim in $\pay(a)$ for any $a \in \mathcal{S}$, or in $y$ after
redaction.
\end{proposition}

\noindent The proposition follows directly from \eqref{eq:dest} and \eqref{eq:flow}, because the broker executes every
sink call and read-only tasks have $G(u) = \emptyset$. It says nothing about $H_\mathit{tamper}$ or about values
transformed by the model, which are left to \rSig{} and the reminder. Section~\ref{sec:discussion} shows that every
remaining bypass either violates one of the three assumptions or falls under $H_\mathit{tamper}$.

\paragraph{Configuration choice and cost.} All choices were made on Dev in pre-declared tuning rounds. \ours{} combines
\cRules{} with \cSand, two existing components, without changing any parameter. This configuration was frozen before
\tOne{} was run and has not changed since. \ours{} calls no additional model. Each check is a substring search over
$u$, $\pay(a)$ or $y$. The reminder lengthens the context by one copy of $u$ after each observation. Latency has not
been measured.
